\chapter{Descent Theory}

Descent theory provides a framework for reconstructing global objects from
local data. Given a covering \(U \to X\), one forms its Cech nerve
\[
U_\bullet \colon
U
\;\substack{\longrightarrow\\[-0.4em]\longrightarrow}\;
U \times_X U
\;\substack{\longrightarrow\\[-0.4em]\longrightarrow\\[-0.4em]\longrightarrow}\;
U \times_X U \times_X U
\;\cdots .
\]
A presheaf-like object \(F\) satisfies descent with respect to this covering
if the value of \(F\) on \(X\) can be recovered from the cosimplicial diagram
\(F(U_\bullet)\). The guiding principle of this chapter is therefore:
\[
F(X) \simeq \Tot\bigl(F(U_\bullet)\bigr).
\]
Here \(\Tot\) should be understood as the appropriate limit of the
cosimplicial descent diagram. In higher-categorical treatments of descent,
this totalization is a homotopy limit over the Cech nerve; for instance,
Khan formulates descent in terms of the totalization of a cosimplicial diagram
associated to an augmented Cech nerve, see \cite{khan2022stacks}.
The purpose of the present chapter is more modest: we only need this principle
in the \(1\)-categorical and \(2\)-categorical cases.

For presheaves of sets, the above condition reduces to the ordinary sheaf
condition. Since \(\Set\) is a \(1\)-category, the totalization is the
usual limit. Thus the descent condition says that \(F(X)\) is the equalizer
\[
F(X)
\longrightarrow
F(U)
\;\substack{\longrightarrow\\[-0.4em]\longrightarrow}\;
F(U \times_X U),
\]
which is precisely the usual gluing condition from
Definition~\ref{def:sheaf-on-cite}. In this case, the higher terms of the
Cech nerve impose no additional coherence beyond equality.

For presheaves valued in groupoids or categories, however, descent data must
remember isomorphisms and their coherences. An object over \(X\) is described
locally by objects over \(U\), together with isomorphisms between their two
pullbacks to \(U \times_X U\). These isomorphisms are not arbitrary: on the
triple overlap \(U \times_X U \times_X U\), they must satisfy the familiar
cocycle condition. In the language of the Cech nerve, this cocycle
condition is the coherence imposed by the \(2\)-simplices of the associated
cosimplicial diagram. Thus the classical descent datum for stacks is exactly
the \(2\)-categorical form of the totalization of \(F(U_\bullet)\).

We shall make this precise without developing the general theory of
homotopy limits. Instead, we use the \(2\)-categorical notion of a
pseudolimit. Gambino shows that, in the setting of \(2\)-categories, homotopy
limits can be studied via model structures on diagram \(2\)-categories and
that pseudolimits are related to such homotopy limits. \cite{gambino2008homotopy}
This justifies taking the pseudolimit of the Cech descent diagram as the
correct \(2\)-categorical replacement for the ordinary limit. Accordingly,
for \(F\) valued in \(\Grpd\) or \(\Cat\), we define descent with respect to
\(U \to X\) by requiring the canonical functor
\[
F(X) \longrightarrow \operatorname{PsLim}_{[n]\in\Delta} F(U_n)
\]
to be an equivalence of groupoids, respectively of categories.

In this way, the chapter treats descent uniformly across the cases needed
below. For set-valued presheaves, descent is expressed by an ordinary limit
and recovers the sheaf equalizer. For groupoid- and category-valued
presheaves, descent is expressed by a pseudolimit and recovers the usual
objects, gluing isomorphisms, and cocycle condition of stack-theoretic
descent.

\section*{Simplicial and cosimplicial objects}

Lets briefly recall the formalism of simplicial and cosimplicial objects
which underlies the Cech nerve and the notion of totalization.

\begin{definition}[Simplicial and cosimplicial objects]
	Let $\Delta$ denote the simplex category, whose objects are the finite ordered sets
	$[n]=\{0,\dots,n\}$ and whose morphisms are non-decreasing maps.
	
	\begin{enumerate}
		\item A simplicial object in a category $\mathcal{C}$ is a functor
		\[
		X_\bullet : \Delta^{op} \to \mathcal{C}.
		\]
		
		\item A cosimplicial object in $\mathcal{C}$ is a functor
		\[
		X^\bullet : \Delta \to \mathcal{C}.
		\]
	\end{enumerate}
	
	Equivalently, a simplicial (resp.\ cosimplicial) object consists of objects
	$X_n$ together with face and degeneracy maps satisfying the simplicial identities.
\end{definition}

\begin{remark}
	A cosimplicial object in $\mathcal{C}$ is the same as a simplicial object
	in the opposite category $\mathcal{C}^{op}$.
	In practice, simplicial objects encode gluing diagrams, while
	cosimplicial objects encode descent data.
\end{remark}

\begin{example}[Cech nerve as a simplicial object]
	Let $f : U \to X$ be a morphism in a category with fiber products.
	The Čech nerve
	\[
	U_\bullet : \Delta^{op} \to \mathcal{C},
	\qquad
	U_n = \underbrace{U \times_X \cdots \times_X U}_{n+1 \text{ times}}.
	\]
	is a simplicial object.
	
	\begin{itemize}
		\item The face maps forget one factor:
		\[
		d^i : U_n \to U_{n-1}, \qquad
		(u_0,\dots,u_n) \mapsto (u_0,\dots,\widehat{u_i},\dots,u_n).
		\]
		
		\item The degeneracy maps repeat a factor via diagonals:
		\[
		s^i : U_n \to U_{n+1}, \qquad
		(u_0,\dots,u_n) \mapsto (u_0,\dots,u_i,u_i,\dots,u_n).
		\]
	\end{itemize}
\end{example}

\begin{example}[From simplicial to cosimplicial via a presheaf]
	Let $F : \mathcal{C}^{op} \to \mathcal{D}$ be a presheaf.
	Applying $F$ to a simplicial object $U_\bullet$ produces a cosimplicial object
	\[
	F(U_\bullet) : \Delta \to \mathcal{D}, \qquad
	[n] \mapsto F(U_n).
	\]
	
	Concretely, the two face maps
	\[
	U \times_X U \rightrightarrows U
	\]
	induce two maps
	\[
	F(U) \rightrightarrows F(U \times_X U),
	\]
	and similarly in higher degrees.
	
	This reversal of arrows is the reason why descent naturally involves
	\emph{cosimplicial} objects.
\end{example}

\begin{example}[Very concrete example: sets]
	Let $F$ be a presheaf of sets and let $U \to X$ be a covering.
	Then the beginning of the cosimplicial diagram $F(U_\bullet)$ is
	\[
	F(U)
	\rightrightarrows
	F(U \times_X U).
	\]
	
	The totalization reduces to the equalizer
	\[
	F(X)
	\cong
	\left\{
	s \in F(U)
	\;\middle|\;
	\pi_1^*(s) = \pi_2^*(s)
	\;\text{in } F(U \times_X U)
	\right\},
	\]
	which is precisely the usual sheaf condition.
\end{example}

TODO: For strict $2$-cateogries
\begin{lemma}[Truncated descent for sets and groupoids]
	\label{lem:truncated-tot-for-stric-n-cats}
	Let $X^\bullet : \Delta \to \mathcal{C}$ be a cosimplicial object in a category $\mathcal{C}$
	
	\begin{enumerate}
		\item If $\mathcal{C} = \Set$, then the natural map
		\[
		\Tot(X^\bullet) \longrightarrow \Tot_1(X^\bullet)
		=
		\Eq\big(X^0 \rightrightarrows X^1\big)
		\]
		is an isomorphism.
		
		\item If $\mathcal{C} = \Grpd$, then the natural map
		\[
		\Tot(X^\bullet) \longrightarrow \Tot_2(X^\bullet)
		\]
		is an equivalence of categories.
	\end{enumerate}
\end{lemma}

\begin{proof}
	We describe explicitly what the totalization computes.
	
	\medskip
	
	(1) In $\Set$, the limit of a cosimplicial diagram consists of elements
	$s \in X^0$ such that the two images in $X^1$ agree.
	Higher compatibility conditions are automatic since equalities are strict.
	Hence $\Tot(X^\bullet)$ is already the equalizer of
	$X^0 \rightrightarrows X^1$.
	
	\medskip
	
	(2) In $\Grpd$, the totalization $\Tot(X^\bullet)$ is the category of descent data consisting of  objects: $(N, \alpha)$, with $N \in X^0$ and an isomorphism
	$\alpha : d^0(N) \to d^1(N)$ in $X^1$ satisfying the cocycle condition
	$d^2(\alpha) \circ d^0(\alpha) = d^1(\alpha)$, i.e. letting the following commute:
	% https://q.uiver.app/#q=WzAsMyxbMCwxLCJkXjBkXjAoTikiXSxbMiwxLCJkXjFkXjEoTikiXSxbMSwwLCJkXjBkXjFOIl0sWzAsMSwiZF4xXFxhbHBoYSIsMl0sWzAsMiwiZF4wXFxhbHBoYSJdLFsyLDEsImReMlxcYWxwaGEiXV0=
	\[\begin{tikzcd}
		& {d^0d^1N} & \\
		{d^0d^0(N)} && {d^1d^1(N)}
		\arrow["{d^2\alpha}", from=1-2, to=2-3]
		\arrow["{d^0\alpha}", from=2-1, to=1-2]
		\arrow["{d^1\alpha}"', from=2-1, to=2-3]
	\end{tikzcd}\]

	Here we applied $d^0, d^1$ and $d^2$ to the morphism $\alpha$ and impliclty used the simplicial identities for the cofaces $d^i$.
	
	All higher compatibility conditions (in degrees $\ge 3$) are automatic,
	because in a groupoid all higher diagrams commute uniquely once the
	2-dimensional compatibility is satisfied.
	
	Hence $\Tot(X^\bullet)$ is already determined by the truncation
	$\Delta_{\le 2}$, i.e.\ $\Tot_2(X^\bullet)$.
\end{proof}

\begin{remark}
	This lemma reflects the general principle that sets are $0$-truncated
	and groupoids are $1$-truncated objects. For $n$-truncated objects,
	all higher coherence conditions above level $n+1$ are automatic, so
	totalizations stabilize at level $n+1$.
\end{remark}

\begin{definition}[Descent datum]
	Let $\CatC$ be a category with pullbacks, $F: \CatC^\op\to \Cat$ a pseudo functor and $U \to X$ a morphism. A descent datum for $F$ along $U \to X$ consists of
	\begin{enumerate}
		\item an object $x\in F(U)$,
		\item an isomorphism $\alpha:d^0(x) \xrightarrow{\sim} d^1(x)$ in $F(U\times_X U))$ between two pullbacks of $x$ to
		$U\times_X U$, such that $s^0(\alpha) = id_\alpha$
		\item a cocycle condition $d^2(\alpha)\circ d^0(\alpha) = d^1(\alpha)$ in $F(U\times_X U\times_X U)$, i.e. the following commutes:
		% https://q.uiver.app/#q=WzAsMyxbMCwxLCJkXjBkXjAoeCkiXSxbMiwxLCJkXjFkXjEoeCkiXSxbMSwwLCJkXjBkXjEoeCkiXSxbMCwxLCJkXjFcXGFscGhhIiwyXSxbMCwyLCJkXjBcXGFscGhhIl0sWzIsMSwiZF4yXFxhbHBoYSJdXQ==
		\[\begin{tikzcd}
			& {d^0d^1(x)} & \\
			{d^0d^0(x)} && {d^1d^1(x)}
			\arrow["{d^2\alpha}", from=1-2, to=2-3]
			\arrow["{d^0\alpha}", from=2-1, to=1-2]
			\arrow["{d^1\alpha}"', from=2-1, to=2-3]
		\end{tikzcd}\]
	\end{enumerate}
	A morphism of descent data
	\[
	f:(x,\alpha)\longrightarrow (y,\beta)
	\]
	is a morphism
	\[
	f:x\to y
	\]
	in $F(U)$ such that the diagram
	\[
	\begin{tikzcd}
		d^0(x) \arrow[r,"\alpha"] \arrow[d,"d^0(f)"'] &
		d^1(x) \arrow[d,"d^1(f)"] \\
		d^0(y) \arrow[r,"\beta"'] &
		d^1(y)
	\end{tikzcd}
	\]
	commutes in $F(U\times_X U)$.
	We write $\Desc_p(F)$ for the category whose objects are descent data and
	whose morphisms are morphisms of descent data.
\end{definition}

\begin{remark}[Descent data as a 2-categorical totalization]
	Let \(p:U\to X\) be a cover and let
	\[
	U_\bullet
	\]
	be its \v{C}ech nerve. Applying the pseudofunctor
	\[
	F:\CatC^{op}\to\Cat
	\]
	gives a pseudocosimplicial category
	\[
	F(U_\bullet):\Delta\to\Cat .
	\]
	From a 2-categorical point of view, the category of descent data should be thought of
	as the pseudo-limit, or descent object, of the beginning of this cosimplicial diagram:
	\[
	\begin{tikzcd}
		{F(U)} & {F(U\times_XU)} & {F(U\times_XU\times_XU)} .
		\arrow[shift left, from=1-1, to=1-2]
		\arrow[shift right, from=1-1, to=1-2]
		\arrow[from=1-2, to=1-3]
		\arrow[shift left=2, from=1-2, to=1-3]
		\arrow[shift right=2, from=1-2, to=1-3]
	\end{tikzcd}
	\]
	Thus an object of this pseudo-limit consists of an object \(x\in F(U)\), together
	with an invertible comparison morphism
	\[
	\alpha:d^0x\xrightarrow{\sim}d^1x
	\]
	in \(F(U\times_XU)\). The compatibility of the pseudo-cone with the degeneracy map
	gives the normalization condition
	\[
	s^0(\alpha)=id_x,
	\]
	and the compatibility with the three coface maps gives the cocycle condition
	\[
	d^2(\alpha)\circ d^0(\alpha)=d^1(\alpha)
	\]
	in \(F(U\times_XU\times_XU)\), up to the coherence isomorphisms of the
	pseudofunctor \(F\). In this sense, the usual cocycle condition is precisely the
	higher coherence condition which appears when the ordinary sheaf condition is
	categorified.
	
	Equivalently, one may write informally
	\[
	\Desc_p(F)
	\simeq
	\operatorname{pslim}_{[n]\in\Delta_{\leq 2}} F(U_n),
	\]
	where the right-hand side denotes the 2-categorical totalization of the
	2-truncated \v{C}ech diagram.
	\footnote{
		Strictly speaking, the displayed pseudo-limit notation suppresses the coherence
		isomorphisms coming from the pseudofunctor \(F\). Hollander gives the corresponding
		explicit description for cosimplicial groupoids: the homotopy limit, denoted
		\(\operatorname{Tot}_2\), has objects \((x,\alpha)\) satisfying a unit condition and
		a cocycle condition; see \cite[Corollary~2.11]{hollander2001homotopy}. In the
		proof of \cite[Theorem~3.9]{hollander2001homotopy}, she identifies such objects
		with classical descent data for a cover. For the \(\Cat\)-valued 2-categorical
		formulation, Nunes defines the descent object of a pseudocosimplicial object as a
		right pseudo-Kan extension \cite[Definition~4.3]{lucatellinunes2018pseudokan}, observes that it
		is a conical bilimit \cite[Remark~4.4]{lucatellinunes2018pseudokan}, compares this with the usual
		weighted definition of the descent object \cite[Definition~4.6 and Theorem~4.11]
		{lucatellinunes2018pseudokan}, and gives an explicit description of strict descent objects in
		\(\Cat\) by an object together with an invertible comparison morphism satisfying
		associativity and identity conditions \cite[Remark~4.23]{lucatellinunes2018pseudokan}.
	}
\end{remark}


\begin{proposition}[Sheaf condition as Cech descent]
	Let $F$ be a sheaf of sets on a site and let $U \to X$ be a covering
	with Cech nerve $U_\bullet$. Then
	\[
	F(X)
	\cong
	\operatorname{Tot}
	\big(
	F(U_\bullet)
	\big).
	\]
	Since sets are $0$-truncated, the totalization reduces to the equalizer
	\[
	F(X)
	\to
	F(U)
	\rightrightarrows
	F(U \times_X U).
	\]
\end{proposition}

\begin{remark}
	Given a simplicial object $U_\bullet$ in a cateogry $\mathcal{C}$. Applying a preseahf $F$ gives rise to a cosimplicial object $F(U_\bullet)$. 
\end{remark}

The Amitsur complex computes the totalization of the Čech nerve of
$\Spec(B) \to \Spec(A)$, and faithfully flat descent for rings states
that this totalization recovers the ring $A$. Thus descent theory in
algebraic geometry ultimately reduces to exactness of the Amitsur complex.

\begin{theorem}[Amitsur complex]
	Let $A \to B$ be a faithfully flat morphism of rings.
	Then the augmented Amitsur complex
	\[
	0 \to A
	\longrightarrow B
	\longrightarrow B \otimes_A B
	\longrightarrow B \otimes_A B \otimes_A B
	\longrightarrow \cdots
	\]
	with differentials given by the alternating sum of the coface maps
	is exact.
\end{theorem}




\begin{theorem}[Faithfully flat descent for rings]\label{thm:faithfully-flat-descent-for-rings}
	Let $A \to B$ be a faithfully flat morphism of rings and let
	% https://q.uiver.app/#q=WzAsNCxbMCwwLCJCIl0sWzEsMCwiQlxcb3RpbWVzX0EgQiJdLFsyLDAsIkJcXG90aW1lc19BQlxcb3RpbWVzX0FCIl0sWzMsMCwiXFxjZG90cyJdLFswLDEsIiIsMCx7Im9mZnNldCI6LTF9XSxbMSwyXSxbMSwyLCIiLDAseyJvZmZzZXQiOi0yfV0sWzEsMiwiIiwwLHsib2Zmc2V0IjoyfV0sWzAsMSwiIiwwLHsib2Zmc2V0IjoxfV0sWzIsMywiIiwwLHsib2Zmc2V0IjotMX1dLFsyLDMsIiIsMCx7Im9mZnNldCI6MX1dLFsyLDMsIiIsMCx7Im9mZnNldCI6M31dLFsyLDMsIiIsMCx7Im9mZnNldCI6LTN9XV0=
	\[
	B_\bullet =
	\big(
	\begin{tikzcd}
		B & {B\otimes_A B} & {B\otimes_AB\otimes_AB} & \cdots
		\arrow[shift left, from=1-1, to=1-2]
		\arrow[shift right, from=1-1, to=1-2]
		\arrow[from=1-2, to=1-3]
		\arrow[shift left=2, from=1-2, to=1-3]
		\arrow[shift right=2, from=1-2, to=1-3]
		\arrow[shift left, from=1-3, to=1-4]
		\arrow[shift right, from=1-3, to=1-4]
		\arrow[shift right=3, from=1-3, to=1-4]
		\arrow[shift left=3, from=1-3, to=1-4]
	\end{tikzcd}
	\big)
	\]
	
	be the Čech nerve of $A \to B$.
	Then the natural map
	\[
	A \longrightarrow \operatorname{Tot}(B^\bullet)
	\]
	is an isomorphism.
\end{theorem}


\begin{proof}
	The cosimplicial ring $B^\bullet$ associated to the morphism
	$A \to B$ is the Amitsur complex of the extension.
	If $A \to B$ is faithfully flat, the augmented Amitsur complex
	\[
	0 \to A \to B \to B \otimes_A B \to B \otimes_A B \otimes_A B \to \cdots
	\]
	is exact. Exactness of this complex implies that $A$ is the limit
	of the cosimplicial diagram $B^\bullet$. Hence
	\[
	A \cong \operatorname{Tot}(B^\bullet).
	\]
\end{proof}

\begin{remark}
	The cosimplicial ring $B^\bullet$ is the Amitsur complex associated to
	$A \to B$. Faithfully flat descent for rings is equivalent to exactness of this
	complex.
\end{remark}

\begin{theorem}[Faithfully flat descent for modules]
	\label{thm:faithfully-flat-descent-modules}
	(\cite[Tag 023N]{stacks-project})
	Let \(R\to S\) be a faithfully flat ring homomorphism, and consider the
	pseudofunctor
	\[
	\Mod_{(-)}:(\Alg_R)^{op}\to\Cat,
	\qquad
	A\longmapsto \Mod_A,
	\]
	which sends a morphism of \(R\)-algebras \(A\to B\) to the extension of
	scalars functor
	\[
	B\otimes_A -:\Mod_A\to\Mod_B.
	\]
	We define
	\[
	\Desc(S/R)
	:=
	\Desc_{R\to S}(\Mod_{(-)}).
	\]
	
	Explicitly, an object of \(\Desc(S/R)\) is an \(S\)-module \(N\), together
	with an isomorphism
	\[
	\alpha:d^0N\xrightarrow{\sim}d^1N
	\]
	in \(\Mod_{S\otimes_R S}\), satisfying
	\[
	s^0(\alpha)=id_N
	\]
	and
	\[
	d^2(\alpha)\circ d^0(\alpha)=d^1(\alpha)
	\]
	in \(\Mod_{S\otimes_R S\otimes_R S}\).
	
	Equivalently, if
	\[
	\pi_1,\pi_2:S\to S\otimes_R S
	\]
	are given by
	\[
	\pi_1(s)=s\otimes 1,
	\qquad
	\pi_2(s)=1\otimes s,
	\]
	then
	\[
	d^0N=\pi_1^*N,
	\qquad
	d^1N=\pi_2^*N,
	\]
	and \(\alpha\) is an isomorphism
	\[
	\alpha:
	\pi_1^*N\xrightarrow{\sim}\pi_2^*N.
	\]
	The cocycle condition becomes
	\[
	\pi_{23}^*(\alpha)\circ \pi_{12}^*(\alpha)
	=
	\pi_{13}^*(\alpha)
	\]
	in \(\Mod_{S\otimes_R S\otimes_R S}\), where
	\[
	\pi_{12},\pi_{13},\pi_{23}:S\otimes_R S
	\to S\otimes_R S\otimes_R S
	\]
	are defined by
	\[
	\pi_{12}(s\otimes t)=s\otimes t\otimes 1,
	\qquad
	\pi_{13}(s\otimes t)=s\otimes 1\otimes t,
	\qquad
	\pi_{23}(s\otimes t)=1\otimes s\otimes t.
	\]
	
	A morphism
	\[
	f:(N,\alpha)\to (N',\alpha')
	\]
	in \(\Desc(S/R)\) is an \(S\)-linear map \(f:N\to N'\) such that
	\[
	\alpha'\circ d^0f=d^1f\circ\alpha.
	\]
	Equivalently, the square
	\[
	\begin{tikzcd}
		\pi_1^*N \arrow[r,"\alpha"] \arrow[d,"\pi_1^*f"'] &
		\pi_2^*N \arrow[d,"\pi_2^*f"] \\
		\pi_1^*N' \arrow[r,"\alpha'"'] &
		\pi_2^*N'
	\end{tikzcd}
	\]
	commutes.
	
	Then the canonical functor
	\[
	\Mod_R\longrightarrow \Desc(S/R)
	\]
	is an equivalence of categories. Explicitly, it sends an \(R\)-module \(M\)
	to
	\[
	M\longmapsto (S\otimes_R M,\alpha_M),
	\]
	where \(\alpha_M\) is the canonical descent isomorphism
	\[
	d^0(S\otimes_R M)
	\xrightarrow{\sim}
	d^1(S\otimes_R M),
	\]
	equivalently
	\[
	\pi_1^*(S\otimes_R M)
	\xrightarrow{\sim}
	\pi_2^*(S\otimes_R M).
	\]
\end{theorem}

\begin{theorem}[Faithfully flat descent for modules]
	\label{thm:faithfully-flat-descent-modules2}
	(\cite[Tag 023N]{stacks-project})
	Let \(R\to S\) be a faithfully flat ring homomorphism. Consider the
	pseudofunctor
	\[
	\Mod_{(-)}:(\Alg_R)^{op}\longrightarrow \Cat,
	\qquad
	A\longmapsto \Mod_A,
	\]
	which sends a morphism of \(R\)-algebras \(A\to B\) to the extension of
	scalars functor
	\[
	B\otimes_A -:\Mod_A\to\Mod_B.
	\]
	
	Set
	\[
	S^n:=\underbrace{S\otimes_R\cdots\otimes_R S}_{n+1\text{ factors}}.
	\]
	The algebras \(S^n\) form the Amitsur cosimplicial \(R\)-algebra
	associated to \(R\to S\). Applying the pseudofunctor \(\Mod_{(-)}\) gives
	a cosimplicial diagram of categories
	\[
	\begin{tikzcd}
		{\Mod_S} &
		{\Mod_{S\otimes_R S}} &
		{\Mod_{S\otimes_R S\otimes_R S}} &
		\cdots .
		\arrow[shift left, from=1-1, to=1-2]
		\arrow[shift right, from=1-1, to=1-2]
		\arrow[shift right=2, from=1-2, to=1-3]
		\arrow[shift left=2, from=1-2, to=1-3]
		\arrow[from=1-2, to=1-3]
		\arrow[from=1-3, to=1-4]
	\end{tikzcd}
	\]
	
	We define
	\[
	\Desc(S/R):=\Desc_{R\to S}(\Mod_{(-)})
	\]
	to be the descent category associated to this pseudofunctor and the
	morphism \(R\to S\) in \(\Alg_R\).
	
	Explicitly, let
	\[
	\pi_1,\pi_2:S\to S\otimes_R S
	\]
	be given by
	\[
	\pi_1(s)=s\otimes 1,
	\qquad
	\pi_2(s)=1\otimes s.
	\]
	With the convention
	\[
	d^0=\pi_1^*,
	\qquad
	d^1=\pi_2^*,
	\]
	an object of \(\Desc(S/R)\) is an \(S\)-module \(N\), together with an
	isomorphism
	\[
	\alpha:d^0N\xrightarrow{\sim}d^1N
	\]
	in \(\Mod_{S\otimes_R S}\). Equivalently,
	\[
	\alpha:
	N\otimes_{S,\pi_1}(S\otimes_R S)
	\xrightarrow{\sim}
	N\otimes_{S,\pi_2}(S\otimes_R S).
	\]
	This isomorphism is required to satisfy the normalization condition
	\[
	s^0(\alpha)=id_N
	\]
	and the cocycle condition
	\[
	d^2(\alpha)\circ d^0(\alpha)=d^1(\alpha)
	\]
	in \(\Mod_{S\otimes_R S\otimes_R S}\).
	
	In tensor notation, this cocycle condition is
	\[
	\pi_{23}^*(\alpha)\circ \pi_{12}^*(\alpha)
	=
	\pi_{13}^*(\alpha),
	\]
	where
	\[
	\pi_{12},\pi_{13},\pi_{23}:S\otimes_R S\to
	S\otimes_R S\otimes_R S
	\]
	are given by
	\[
	\pi_{12}(s\otimes t)=s\otimes t\otimes 1,
	\]
	\[
	\pi_{13}(s\otimes t)=s\otimes 1\otimes t,
	\]
	\[
	\pi_{23}(s\otimes t)=1\otimes s\otimes t.
	\]
	
	A morphism
	\[
	f:(N,\alpha)\to (N',\alpha')
	\]
	in \(\Desc(S/R)\) is an \(S\)-linear morphism \(f:N\to N'\) such that the
	square
	\[
	\begin{tikzcd}
		d^0N \arrow[r,"\alpha"] \arrow[d,"d^0f"'] &
		d^1N \arrow[d,"d^1f"] \\
		d^0N' \arrow[r,"\alpha'"'] &
		d^1N'
	\end{tikzcd}
	\]
	commutes. Equivalently,
	\[
	\alpha'\circ d^0f=d^1f\circ\alpha.
	\]
	
	Then the canonical functor
	\[
	\Mod_R\longrightarrow \Desc(S/R)
	\]
	is an equivalence of categories. Explicitly, it sends an \(R\)-module \(M\)
	to
	\[
	M\longmapsto (S\otimes_R M,\alpha_M),
	\]
	where \(\alpha_M\) is the canonical isomorphism
	\[
	d^0(S\otimes_R M)
	\xrightarrow{\sim}
	d^1(S\otimes_R M).
	\]
\end{theorem}


\begin{proof}
	For an $R$-module $M$, set
	\[
	F(M):=(S\otimes_R M,\alpha_M),
	\]
	where $\alpha_M$ is the canonical descent datum. This defines a functor
	\[
	F:\Mod_R\to \Desc(S/R).
	\]
	
	Conversely, let $(N,\alpha)\in \Desc(S/R)$. Define
	\[
	G(N,\alpha):=
	\{\,n\in N \mid \alpha(1\otimes n)=n\otimes 1\,\}.
	\]
	Equivalently, $G(N,\alpha)$ is the equalizer of the two maps
	\[
	N \rightrightarrows N\otimes_R S,
	\qquad
	n\mapsto \alpha(n\otimes 1),\quad n\mapsto n\otimes 1.
	\]
	This gives a functor
	\[
	G:\Desc(S/R)\to \Mod_R.
	\]
	
	We show that $F$ and $G$ are quasi-inverse.
	
	First, let $M\in \Mod_R$. Then
	\[
	G(F(M))
	=
	\{\,x\in S\otimes_R M \mid x\otimes 1 = \alpha_M(1\otimes x)\,\}.
	\]
	By the exactness of the Amitsur complex for the faithfully flat map
	$R\to S$, the sequence
	\[
	0\to M \to S\otimes_R M \rightrightarrows S\otimes_R S\otimes_R M
	\]
	is exact. Hence the equalizer above is canonically identified with $M$.
	Thus
	\[
	G(F(M))\cong M.
	\]
	
	Now let $(N,\alpha)\in \Desc(S/R)$. Set $M:=G(N,\alpha)$. There is a natural
	$S$-linear map
	\[
	\theta:S\otimes_R M\longrightarrow N,\qquad s\otimes m\longmapsto sm.
	\]
	We claim that $\theta$ is an isomorphism.
	
	To prove this, consider the augmented twisted Amitsur complex $s(N_\bullet)$ associated to the descent datum $(N,\alpha)$, given by $s(N_n) = N \otimes_R S^{\otimes_R n}$ with $s(N_{-1}) = M$:
	% https://q.uiver.app/#q=WzAsNixbMCwwLCIwIl0sWzEsMCwiTSJdLFsyLDAsIk4iXSxbMywwLCJOXFxvdGltZXNfUiBTIl0sWzQsMCwiTlxcb3RpbWVzX1IgU1xcb3RpbWVzX1IgUyJdLFs1LDAsIlxcY2RvdHMiXSxbMCwxXSxbMSwyLCJcXGlvdGEiXSxbMiwzLCJuXFxtYXBzdG8gXFxhbHBoYSgxXFxvdGltZXMgbikiLDIseyJvZmZzZXQiOjF9XSxbMiwzLCJuXFxtYXBzdG8gblxcb3RpbWVzIDEiLDAseyJvZmZzZXQiOi0xfV0sWzMsNF0sWzMsNCwiIiwwLHsib2Zmc2V0IjotMn1dLFszLDQsIiIsMCx7Im9mZnNldCI6Mn1dLFs0LDUsIiIsMCx7Im9mZnNldCI6LTF9XSxbNCw1LCIiLDAseyJvZmZzZXQiOi0zfV0sWzQsNSwiIiwwLHsib2Zmc2V0IjoxfV0sWzQsNSwiIiwwLHsib2Zmc2V0IjozfV1d
	\[\begin{tikzcd}
		0 & M & N & {N\otimes_R S} & {N\otimes_R S\otimes_R S} & \cdots
		\arrow[from=1-1, to=1-2]
		\arrow["\iota", from=1-2, to=1-3]
		\arrow["{n\mapsto \alpha(1\otimes n)}"', shift right, from=1-3, to=1-4]
		\arrow["{n\mapsto n\otimes 1}", shift left, from=1-3, to=1-4]
		\arrow[from=1-4, to=1-5]
		\arrow[shift left=2, from=1-4, to=1-5]
		\arrow[shift right=2, from=1-4, to=1-5]
		\arrow[shift left, from=1-5, to=1-6]
		\arrow[shift left=3, from=1-5, to=1-6]
		\arrow[shift right, from=1-5, to=1-6]
		\arrow[shift right=3, from=1-5, to=1-6]
	\end{tikzcd}\]
	This is exact by construction of $M$ and faithfully flatness of $R\to S$, i.e. $H^0(s(N_\bullet)) = M$.
	
	Now we tensor this complex with $S\otimes_R - $ and untwist the sequence using the descent datum $\alpha$:
	
	% https://q.uiver.app/#q=WzAsMTAsWzAsMCwiMCJdLFsxLDAsIlNcXG90aW1lc19SIE0iXSxbMiwwLCJTXFxvdGltZXNfUk4iXSxbMywwLCJTXFxvdGltZXNfUiBOXFxvdGltZXNfUiBTIl0sWzQsMCwiU1xcb3RpbWVzX1IgTlxcb3RpbWVzX1IgU1xcb3RpbWVzX1IgUyJdLFswLDEsIjAiXSxbMSwxLCJTXFxvdGltZXNfUk0iXSxbMiwxLCJOXFxvdGltZXNfUiBTIl0sWzMsMSwiTlxcb3RpbWVzX1IgU1xcb3RpbWVzX1IgUyJdLFs0LDEsIk5cXG90aW1lc19SIFNcXG90aW1lc19SIFNcXG90aW1lc19SIFMiXSxbMCwxXSxbMSwyXSxbMiwzLCIiLDAseyJvZmZzZXQiOi0xfV0sWzIsMywiIiwwLHsib2Zmc2V0IjoxfV0sWzMsNF0sWzMsNCwiIiwwLHsib2Zmc2V0IjotMn1dLFszLDQsIiIsMCx7Im9mZnNldCI6Mn1dLFsxLDYsIj0iXSxbNSw2XSxbNiw3LCJcXHBzaSJdLFsyLDcsIlxcYWxwaGEiXSxbNyw4LCIiLDAseyJvZmZzZXQiOi0xfV0sWzMsOCwiXFxhbHBoYVxcb3RpbWVzIGlkIl0sWzgsOV0sWzgsOSwiIiwwLHsib2Zmc2V0IjotMn1dLFs4LDksIiIsMCx7Im9mZnNldCI6Mn1dLFs0LDksIlxcYWxwaGFcXG90aW1lcyBpZFxcb3RpbWVzIGlkIl0sWzcsOCwiIiwwLHsib2Zmc2V0IjoxfV1d
	\[\begin{tikzcd}
		0 & {S\otimes_R M} & {S\otimes_RN} & {S\otimes_R N\otimes_R S} & {S\otimes_R N\otimes_R S\otimes_R S} \\
		0 & {S\otimes_RM} & {N\otimes_R S} & {N\otimes_R S\otimes_R S} & {N\otimes_R S\otimes_R S\otimes_R S}
		\arrow[from=1-1, to=1-2]
		\arrow[from=1-2, to=1-3]
		\arrow["{=}", from=1-2, to=2-2]
		\arrow[shift left, from=1-3, to=1-4]
		\arrow[shift right, from=1-3, to=1-4]
		\arrow["\alpha", from=1-3, to=2-3]
		\arrow[from=1-4, to=1-5]
		\arrow[shift left=2, from=1-4, to=1-5]
		\arrow[shift right=2, from=1-4, to=1-5]
		\arrow["{\alpha\otimes id}", from=1-4, to=2-4]
		\arrow["{\alpha\otimes id\otimes id}", from=1-5, to=2-5]
		\arrow[from=2-1, to=2-2]
		\arrow["\psi", from=2-2, to=2-3]
		\arrow[shift left, from=2-3, to=2-4]
		\arrow[shift right, from=2-3, to=2-4]
		\arrow[from=2-4, to=2-5]
		\arrow[shift left=2, from=2-4, to=2-5]
		\arrow[shift right=2, from=2-4, to=2-5]
	\end{tikzcd}\]

	The morphism $\psi: S\otimes_R M \to N\otimes_R S$ factors as 
	% https://q.uiver.app/#q=WzAsMyxbMCwwLCJTXFxvdGltZXNfUk0iXSxbMSwwLCJTXFxvdGltZXNfUk4iXSxbMiwwLCJOXFxvdGltZXNfUlMiXSxbMCwxLCJTXFxvdGltZXNfUlxcaW90YSJdLFsxLDIsIlxcYWxwaGEiXV0=
	\[\begin{tikzcd}
		{S\otimes_RM} & {S\otimes_RN} & {N\otimes_RS},
		\arrow["{S\otimes_R\iota}", from=1-1, to=1-2]
		\arrow["\alpha", from=1-2, to=1-3]
	\end{tikzcd}\]

	where $\iota$ is the inclusion $M\subset N$. Since for all $m\in M$ we have $\alpha(1\otimes m) = m\otimes 1$, we get 
	$\psi(s\otimes m) = \alpha(s\otimes m) = \alpha((s\otimes 1)(1\otimes m)) = (s\otimes 1)\alpha(1\otimes m) = (s\otimes 1)(m\otimes 1) = (sm)\otimes 1$ by $S\otimes_R S$-linearity of $\alpha$. Hence, this morphisms $\psi$ is given by $s\otimes m\mapsto (sm)\otimes 1$.
	
	Now we compare the second row to the Amitsur complex $N_\bullet$ of $N$:
	% https://q.uiver.app/#q=WzAsMTUsWzAsMCwiMCJdLFsxLDAsIlNcXG90aW1lc19SIE0iXSxbMiwwLCJTXFxvdGltZXNfUk4iXSxbMywwLCJTXFxvdGltZXNfUiBOXFxvdGltZXNfUiBTIl0sWzQsMCwiU1xcb3RpbWVzX1IgTlxcb3RpbWVzX1IgU1xcb3RpbWVzX1IgUyJdLFswLDEsIjAiXSxbMSwxLCJTXFxvdGltZXNfUk0iXSxbMiwxLCJOXFxvdGltZXNfUiBTIl0sWzMsMSwiTlxcb3RpbWVzX1IgU1xcb3RpbWVzX1IgUyJdLFs0LDEsIk5cXG90aW1lc19SIFNcXG90aW1lc19SIFNcXG90aW1lc19SIFMiXSxbMCwyLCIwIl0sWzEsMiwiTiJdLFsyLDIsIk5cXG90aW1lc19SIFMiXSxbMywyLCJOXFxvdGltZXNfUiBTXFxvdGltZXNfUiBTIl0sWzQsMiwiTlxcb3RpbWVzX1IgU1xcb3RpbWVzX1IgU1xcb3RpbWVzX1IgUyJdLFswLDFdLFsxLDJdLFsyLDMsIiIsMCx7Im9mZnNldCI6LTF9XSxbMiwzLCIiLDAseyJvZmZzZXQiOjF9XSxbMyw0XSxbMyw0LCIiLDAseyJvZmZzZXQiOi0yfV0sWzMsNCwiIiwwLHsib2Zmc2V0IjoyfV0sWzEsNiwiPSJdLFs1LDZdLFs2LDcsIlxccHNpIl0sWzIsNywiXFxhbHBoYSJdLFs3LDgsIiIsMCx7Im9mZnNldCI6LTF9XSxbMyw4LCJcXGFscGhhXFxvdGltZXMgaWQiXSxbOCw5XSxbOCw5LCIiLDAseyJvZmZzZXQiOi0yfV0sWzgsOSwiIiwwLHsib2Zmc2V0IjoyfV0sWzQsOSwiXFxhbHBoYVxcb3RpbWVzIGlkXFxvdGltZXMgaWQiXSxbNyw4LCIiLDAseyJvZmZzZXQiOjF9XSxbNiwxMSwiXFx0aGV0YSJdLFsxMCwxMV0sWzExLDEyXSxbMTIsMTMsIiIsMix7Im9mZnNldCI6MX1dLFsxMiwxMywiIiwyLHsib2Zmc2V0IjotMX1dLFsxMywxNF0sWzEzLDE0LCIiLDIseyJvZmZzZXQiOi0yfV0sWzEzLDE0LCIiLDIseyJvZmZzZXQiOjJ9XSxbNywxMiwiPSJdLFs4LDEzLCI9Il0sWzksMTQsIj0iXV0=
	\[\begin{tikzcd}
		0 & {S\otimes_R M} & {S\otimes_RN} & {S\otimes_R N\otimes_R S} & {S\otimes_R N\otimes_R S\otimes_R S} \\
		0 & {S\otimes_RM} & {N\otimes_R S} & {N\otimes_R S\otimes_R S} & {N\otimes_R S\otimes_R S\otimes_R S} \\
		0 & N & {N\otimes_R S} & {N\otimes_R S\otimes_R S} & {N\otimes_R S\otimes_R S\otimes_R S}
		\arrow[from=1-1, to=1-2]
		\arrow[from=1-2, to=1-3]
		\arrow["{=}", from=1-2, to=2-2]
		\arrow[shift left, from=1-3, to=1-4]
		\arrow[shift right, from=1-3, to=1-4]
		\arrow["\alpha", from=1-3, to=2-3]
		\arrow[from=1-4, to=1-5]
		\arrow[shift left=2, from=1-4, to=1-5]
		\arrow[shift right=2, from=1-4, to=1-5]
		\arrow["{\alpha\otimes id}", from=1-4, to=2-4]
		\arrow["{\alpha\otimes id\otimes id}", from=1-5, to=2-5]
		\arrow[from=2-1, to=2-2]
		\arrow["\psi", from=2-2, to=2-3]
		\arrow["\theta", from=2-2, to=3-2]
		\arrow[shift left, from=2-3, to=2-4]
		\arrow[shift right, from=2-3, to=2-4]
		\arrow["{=}", from=2-3, to=3-3]
		\arrow[from=2-4, to=2-5]
		\arrow[shift left=2, from=2-4, to=2-5]
		\arrow[shift right=2, from=2-4, to=2-5]
		\arrow["{=}", from=2-4, to=3-4]
		\arrow["{=}", from=2-5, to=3-5]
		\arrow[from=3-1, to=3-2]
		\arrow[from=3-2, to=3-3]
		\arrow[shift right, from=3-3, to=3-4]
		\arrow[shift left, from=3-3, to=3-4]
		\arrow[from=3-4, to=3-5]
		\arrow[shift left=2, from=3-4, to=3-5]
		\arrow[shift right=2, from=3-4, to=3-5]
	\end{tikzcd}\]
	
	By definition of $\theta$, this diagram commutes, i.e. $\psi(s\otimes m) = sm\otimes 1 = \theta(s\otimes m)\otimes 1$.
	
	Since $R\to S$ is faithfully flat, the Amitsur $N_\bullet$ complex is exact.
	
	% https://q.uiver.app/#q=WzAsMTAsWzAsMCwiMCJdLFsxLDAsIlNcXG90aW1lc19STSJdLFsyLDAsIlNcXG90aW1lc19SIE4iXSxbMywwLCJTXFxvdGltZXNfUiBOXFxvdGltZXNfUiBTIl0sWzQsMCwiXFxjZG90cyJdLFswLDEsIjAiXSxbMSwxLCJOIl0sWzIsMSwiTlxcb3RpbWVzX1IgUyJdLFszLDEsIk5cXG90aW1lc19SIFNcXG90aW1lc19SIFMiXSxbNCwxLCJcXGNkb3RzIl0sWzAsMV0sWzEsMiwiXFxwc2kiXSxbMiwzLCIiLDAseyJvZmZzZXQiOi0xfV0sWzMsNF0sWzMsNCwiIiwwLHsib2Zmc2V0IjotMn1dLFszLDQsIiIsMCx7Im9mZnNldCI6Mn1dLFsyLDMsIiIsMCx7Im9mZnNldCI6MX1dLFsxLDYsIlxcdGhldGEiXSxbNSw2XSxbNiw3XSxbNyw4LCIiLDIseyJvZmZzZXQiOjF9XSxbNyw4LCIiLDIseyJvZmZzZXQiOi0xfV0sWzgsOV0sWzgsOSwiIiwyLHsib2Zmc2V0IjotMn1dLFs4LDksIiIsMix7Im9mZnNldCI6Mn1dLFsyLDcsIj0iXSxbMyw4LCI9Il1d
	\[\begin{tikzcd}
		0 & {S\otimes_RM} & {S\otimes_R N} & {S\otimes_R N\otimes_R S} & \cdots \\
		0 & N & {N\otimes_R S} & {N\otimes_R S\otimes_R S} & \cdots
		\arrow[from=1-1, to=1-2]
		\arrow["\psi", from=1-2, to=1-3]
		\arrow["\theta", from=1-2, to=2-2]
		\arrow[shift left, from=1-3, to=1-4]
		\arrow[shift right, from=1-3, to=1-4]
		\arrow["{=}", from=1-3, to=2-3]
		\arrow[from=1-4, to=1-5]
		\arrow[shift left=2, from=1-4, to=1-5]
		\arrow[shift right=2, from=1-4, to=1-5]
		\arrow["{=}", from=1-4, to=2-4]
		\arrow[from=2-1, to=2-2]
		\arrow[from=2-2, to=2-3]
		\arrow[shift right, from=2-3, to=2-4]
		\arrow[shift left, from=2-3, to=2-4]
		\arrow[from=2-4, to=2-5]
		\arrow[shift left=2, from=2-4, to=2-5]
		\arrow[shift right=2, from=2-4, to=2-5]
	\end{tikzcd}\]
	
	
	
	 Hence, we obtain an isomorphism
	\[
	S\otimes_R M \xrightarrow{\theta} H^0(N_\bullet) \cong N
	\]
	Moreover, under this isomorphism the canonical descent
	datum on $S\otimes_R M$ identifies with the given descent datum $\alpha$ on
	$N$. Hence
	\[
	F(G(N,\alpha))\cong (N,\alpha).
	\]
	
	Therefore $F$ and $G$ are quasi-inverse equivalences.
\end{proof}

\begin{definition}[Descent category for quasi-coherent sheaves]
	Let \(f:U\to X\) be a morphism of schemes. Consider the pseudofunctor
	\[
	\QCoh:\Sch^{op}\to\Cat,
	\qquad
	X\longmapsto \QCoh(X).
	\]
	We write
	\[
	\Desc_{\QCoh}(U/X)
	:=
	\Desc_f(\QCoh)
	\]
	for the category of descent data for \(\QCoh\) along \(f\).
	
	Explicitly, an object of \(\Desc_{\QCoh}(U/X)\) is a pair
	\[
	(\mathcal F,\varphi)
	\]
	consisting of a quasi-coherent sheaf \(\mathcal F\in\QCoh(U)\), together with
	an isomorphism
	\[
	\varphi:\pi_1^*\mathcal F\xrightarrow{\sim}\pi_2^*\mathcal F
	\]
	in \(\QCoh(U\times_X U)\), where
	\[
	\pi_1,\pi_2:U\times_X U\to U
	\]
	are the two projections.
	
	This isomorphism is required to satisfy the normalization condition
	\[
	\Delta^*(\varphi)=id_{\mathcal F},
	\]
	where
	\[
	\Delta:U\to U\times_X U
	\]
	is the diagonal, and the cocycle condition
	\[
	\pi_{23}^*(\varphi)\circ \pi_{12}^*(\varphi)
	=
	\pi_{13}^*(\varphi)
	\]
	in \(\QCoh(U\times_X U\times_X U)\). Here
	\[
	\pi_{ij}:U\times_X U\times_X U\to U\times_X U
	\]
	denotes the projection onto the \(i\)-th and \(j\)-th factors.
	
	A morphism
	\[
	\psi:(\mathcal F,\varphi)\longrightarrow(\mathcal F',\varphi')
	\]
	is a morphism
	\[
	\psi:\mathcal F\to\mathcal F'
	\]
	in \(\QCoh(U)\) such that the diagram
	\[
	\begin{tikzcd}
		\pi_1^*\mathcal F
		\arrow[r,"\varphi"]
		\arrow[d,"\pi_1^*\psi"']
		&
		\pi_2^*\mathcal F
		\arrow[d,"\pi_2^*\psi"]
		\\
		\pi_1^*\mathcal F'
		\arrow[r,"\varphi'"']
		&
		\pi_2^*\mathcal F'
	\end{tikzcd}
	\]
	commutes in \(\QCoh(U\times_X U)\). Equivalently,
	\[
	\varphi'\circ \pi_1^*(\psi)
	=
	\pi_2^*(\psi)\circ \varphi .
	\]
\end{definition}

\begin{theorem}[Faithfully flat descent for quasi-coherent sheaves]
	(\cite[Tag 023T]{stacks-project})
	Let $f : U \to X$ be an fpqc covering of schemes and let $U_\bullet$
	be the Čech nerve.
	
	Then the pullback functor induces an equivalence of categories
	\[
	\QCoh(X)
	\longrightarrow
	\Desc_{\QCoh}(U/X),
	\]
	where $\Desc_{\QCoh}(U/X)$ denotes the category of descent data
	for quasi-coherent sheaves along $U \to X$.
\end{theorem}

\begin{proof}
	We first reduce to the affine case and then glue the local statements.
	
	Choose an affine open basis
	\[
	X = \bigcup_{j \in J} V_j,
	\qquad
	V_j = \Spec A_j,
	\]
	which is closed under intersection.
	Then each intersection $V_j\cap V_{j'}$ is again affine.
	For each $j$, set
	\[
	U_j := U \times_X V_j.
	\]
	Since fpqc morphisms are stable under base change, each morphism
	\[
	U_j \to V_j
	\]
	is again an fpqc covering.
	
	We claim that it suffices to prove that for each $j$ the pullback functor
	\begin{align}\label{diag:descent-qcoh-affine-restriction}
			\QCoh(V_j) \longrightarrow \Desc_{\QCoh} (U_j/X)
	\end{align}
	is an equivalence.
	Assume \ref{diag:descent-qcoh-affine-restriction} is an equivalence for all $j$.
	For essential surjectivity, let
	\[
	(\mathcal F_U,\alpha)\in \Desc_{\QCoh}(U/X)
	\]
	be a descent datum on $U$. Restricting to $U_j$ yields for each $j$ an object
	\[
	(\mathcal F_{U_j},\alpha_j)\in  \Desc_{\QCoh} (U_j/X)).
	\]
	Assuming the affine case, there exists a quasi-coherent sheaf
	\[
	\mathcal F_j\in \QCoh(V_j)
	\]
	whose pullback along $U_j\to V_j$ is identified with
	$(\mathcal F_{U_j},\alpha_j)$.
	Write $U_{jj'}$ for the pullback $U\times_X(V_j\cap V_{j'})$.
	In the following diagram, the squiggly morphisms denote the category over the schemes $U,U_J, V_j, X$.
	
	% https://q.uiver.app/#q=WzAsMTIsWzAsNCwiVSJdLFsxLDQsIlgiXSxbMiwzLCJWX2oiXSxbMSwzLCJVX2oiXSxbMCwyLCJcXERlc2Nfe1xcUUNvaH0oVS9YKSJdLFsxLDEsIlxcRGVzY197XFxRQ29ofShVX2opIl0sWzIsMSwiXFxRQ29oKFZfaikiXSxbMSwyLCJcXFFDb2goWCkiXSxbMywyLCJWX2pcXGNhcCBWX3tqJ30iXSxbMiwyLCJVX3tqaid9Il0sWzIsMCwiXFxEZXNjX3tcXFFDb2h9KFVfe2pqJ30pIl0sWzMsMCwiXFxRQ29oKFZfalxcY2FwIFZfe2onfSkiXSxbMCwxXSxbMiwxXSxbMywwXSxbMywyXSxbNCwwLCIiLDIseyJzdHlsZSI6eyJib2R5Ijp7Im5hbWUiOiJzcXVpZ2dseSJ9fX1dLFs1LDMsIiIsMix7InN0eWxlIjp7ImJvZHkiOnsibmFtZSI6InNxdWlnZ2x5In19fV0sWzMsMSwiIiwyLHsic3R5bGUiOnsibmFtZSI6ImNvcm5lciJ9fV0sWzYsMiwiIiwyLHsib2Zmc2V0IjotMywic3R5bGUiOnsiYm9keSI6eyJuYW1lIjoic3F1aWdnbHkifX19XSxbNSw2LCJcXGNvbmciXSxbNywxLCIiLDAseyJvZmZzZXQiOi0zLCJzdHlsZSI6eyJib2R5Ijp7Im5hbWUiOiJzcXVpZ2dseSJ9fX1dLFs0LDddLFs2LDddLFs1LDRdLFs5LDhdLFs5LDNdLFs4LDJdLFsxMCw5LCIiLDAseyJzdHlsZSI6eyJib2R5Ijp7Im5hbWUiOiJzcXVpZ2dseSJ9fX1dLFsxMSw4LCIiLDIseyJzdHlsZSI6eyJib2R5Ijp7Im5hbWUiOiJzcXVpZ2dseSJ9fX1dLFsxMCw1XSxbMTEsNl0sWzEwLDExLCJcXGNvbmciXV0=
	\[\begin{tikzcd}
		&& {\Desc_{\QCoh}(U_{jj'})} & {\QCoh(V_j\cap V_{j'})} \\
		& {\Desc_{\QCoh}(U_j)} & {\QCoh(V_j)} \\
		{\Desc_{\QCoh}(U/X)} & {\QCoh(X)} & {U_{jj'}} & {V_j\cap V_{j'}} \\
		& {U_j} & {V_j} \\
		U & X
		\arrow["\cong", from=1-3, to=1-4]
		\arrow[from=1-3, to=2-2]
		\arrow[squiggly, from=1-3, to=3-3]
		\arrow[from=1-4, to=2-3]
		\arrow[squiggly, from=1-4, to=3-4]
		\arrow["\cong", from=2-2, to=2-3]
		\arrow[from=2-2, to=3-1]
		\arrow[squiggly, from=2-2, to=4-2]
		\arrow[from=2-3, to=3-2]
		\arrow[shift left=3, squiggly, from=2-3, to=4-3]
		\arrow[from=3-1, to=3-2]
		\arrow[squiggly, from=3-1, to=5-1]
		\arrow[shift left=3, squiggly, from=3-2, to=5-2]
		\arrow[from=3-3, to=3-4]
		\arrow[from=3-3, to=4-2]
		\arrow[from=3-4, to=4-3]
		\arrow[from=4-2, to=4-3]
		\arrow[from=4-2, to=5-1]
		\arrow["\lrcorner"{anchor=center, pos=0.125, rotate=-45}, draw=none, from=4-2, to=5-2]
		\arrow[from=4-3, to=5-2]
		\arrow[from=5-1, to=5-2]
	\end{tikzcd}\]

	By assumption on the cover, the intersection $V_j\cap V_{j'}$ is affine.
	Now for $j,j'\in J$, the restrictions
	\[
	\mathcal F_j|_{V_j\cap V_{j'}}
	\qquad\text{and}\qquad
	\mathcal F_{j'}|_{V_j\cap V_{j'}}
	\]
	induce the same descent datum on
	\[
	U_{jj'} \to V_j\cap V_{j'},
	\]
	because both $\mathcal{F}_j|_{V_j\cap V_{j'}}$ and $\mathcal{F}_{j'}|_{V_j\cap V_{j'}}$ pull back to the same object $\mathcal{F}_U|_{U_{jj'}}$ with the same cocycle isomorphism $\alpha$.
	Hence they define the same object of 
	\[
	\Desc_{\QCoh}(U_{jj'}).
	\]
	
	By the affine case, for $V_j\cap V_{j'}$ and the fpqc morphism 
	\[
	U_{jj'} \to V_j\cap V_{j'},
	\]
	there is therefore a unique isomorphism
	\[
	\theta_{jj'}:
	\mathcal F_j|_{V_j\cap V_{j'}}
	\xrightarrow{\sim}
	\mathcal F_{j'}|_{V_j\cap V_{j'}}
	\]
	whose pullback to $U_{jj'}$ is the identity on $\mathcal{F}_U|_{U_{jj'}}$. By uniqueness, these isomorphisms satisfy
	the cocycle condition on triple overlaps. Hence the sheaves $\mathcal F_j$
	glue to a quasi-coherent sheaf
	\[
	\mathcal F\in \QCoh(X).
	\]
	By construction, $f^*\mathcal{F}$ with its canonical descent datum is identified with $(\mathcal{F}_U, \alpha)$.
	This proves essential surjectivity globally.
	
	Now for fully faithfulness let $\mathcal{F}, \mathcal{G}\in \QCoh(X)$. We show that there is a bijection
	\[
	\Hom_{\QCohX}(\mathcal{F}, \mathcal{G}) \to \Hom_{\Desc_{\QCoh}(U/X)}(f^*\mathcal{F}, f^*\mathcal{G}).
	\]	
	
	Injectivity: 
	
	Let $\varphi, \psi: \mathcal{F}\to \mathcal{G}$ have the same pullback to $U$, then for every $V_j$ their restrictions
	\[
	\varphi|_{V_j}, \psi|_{V_j}: \mathcal{F}|_{V_j} \to \mathcal{G}|_{V_j}
	\]
	have the same pullback to $U_j$. By the affine case for $V_j$, we get $\varphi|_{V_j} = \psi|_{V_j}$. Since they agree on an open cover, $\varphi = \psi$. 
	
	Surjectivity:
	
	Let $\theta:f^*\mathcal{F} \to f^*\mathcal{G}$ be a morphism in $\Desc_{\QCoh}(U/X)$. Restricting to $U_j$ yields a morphism in $\theta_{j}:f^*\mathcal{F}|_{U_j}\to f^*\mathcal{G}|_{U_j}$ in $\Desc_{\QCoh}(U_j/X)$. By the affine case this corresponds to a morphism $\mathcal{F}|_{U_j} \to \mathcal{G}|_{U_j}$. Now for $j, j'$, the morphism
	$\theta_{j}|_{U_{jj'}}$ and $\theta_{j'}|_{U_{jj'}}$ restrict to the same pullback $\theta_{jj'}$ in 
	\[
	\Desc_{\QCoh}(U_{jj'}/X),
	\]
	which corresponds to a morphism $\mathcal{F}|_{U_{jj'}} \to \mathcal{G}|_{U_{jj'}}$, again by the affine version. 
	By the sheaf property, these morphisms agree on overlaps and therefore glue uniquely to a morphism
	\[
	\mathcal F \to \mathcal G.
	\]
	By construction, this morphism pulls back to $\theta$ in $\Desc_{\QCoh}(U/X)$, which proves surjectivity.
	
	Hence, it remains to prove the affine case. Thus assume
	\[
	X=\Spec A
	\]
	is affine. Since $U\to X$ is fpqc, we may choose a finite affine open covering
	\[
	U=\bigcup_{i\in I} U_i,
	\qquad
	U_i=\Spec B_i.
	\]
	Then
	\[
	\coprod_{i\in I} U_i \to X
	\]
	is again an fpqc covering, and its Cech nerve refines the Cech nerve of
	$U\to X$.
	Since descent for quasi-coherent sheaves is local on the source, replacing $U$ by the fpqc covering $\coprod_{i\in I} U_i \to X$ does not change the category of descent data (\cite[Tag 02L4]{stacks-project}). Hence, it suffices to prove the statement for the affine fpqc cover
	\[
	\coprod_{i\in I}\Spec B_i \to \Spec A.
	\]
	For an affine scheme, we have equivalence of categories, see \cite[Tag 01IB]{stacks-project}:
	\[
	\QCoh(\Spec A)\cong \Mod_A
	\]
	and
	\[
	\QCoh(\Spec B_i)\cong \Mod_{B_i}
	\]
	Under these identifications, the Cech nerve of the covering
	\[
	\coprod_i \Spec B_i \to \Spec A
	\]
	corresponds to the cosimplicial ring obtained from the fpqc algebra
	\[
	A \to \prod_i B_i.
	\]
	Since a finite product of faithfully flat $A$-algebras is again faithfully
	flat, descent for quasi-coherent sheaves over this affine fpqc cover reduces
	to faithfully flat descent for modules (Thm \ref{thm:faithfully-flat-descent-modules}):
	\[
	\Mod_A \cong \Desc(B/A)),
	\]
	where $B=\prod_i B_i$ and $B^\bullet$ is the associated Amitsur-Cech nerve.
\end{proof}


\begin{remark}
	This shows that quasi-coherent sheaves satisfy effective descent with respect
	to fpqc coverings, hence define a stack in groupoids on the fpqc site.
\end{remark}


\begin{theorem}[Descent for affine schemes via relative spectrum]
	Let $X$ be a scheme and let $U \to X$ be an fpqc covering with Cech nerve $U_\bullet$.
	Then the functor
	\[
	\Aff/X \;\longrightarrow\; \Tot\big((\Aff/U_\bullet)\big)
	\]
	is an equivalence, where $\Aff/X$ denotes the category of affine schemes over $X$.
	
	Equivalently, for quasi-coherent $\mathcal{O}_X$-algebras one has
	\[
	\QCohAlg(X)
	\cong
	\Tot\big(\QCohAlg(U_\bullet)\big),
	\]
	and this equivalence is compatible with the relative spectrum construction
	\[
	\underline{\Spec}_X : \QCohAlg(X)^{op} \;\longrightarrow\; \Aff/X.
	\]
\end{theorem}

\begin{proof}
	We first treat quasi-coherent algebras and then pass to affine schemes via
	relative $\Spec$.
	
	Since quasi-coherent $\mathcal{O}_X$-modules satisfy fpqc descent, the same
	is true for quasi-coherent $\mathcal{O}_X$-algebras: the algebra structure
	(consisting of multiplication and unit maps) is given by morphisms in
	$\QCoh(X)$ and is preserved under limits. Hence
	\[
	\QCohAlg(X)
	\cong
	\Tot\big(\QCohAlg(U_\bullet)\big).
	\]
	
	Now consider the relative spectrum construction
	\[
	\underline{\Spec}_X : \QCohAlg(X)^{op} \to \Aff/X.
	\]
	This functor is functorial and compatible with base change. Moreover, it
	induces an equivalence between quasi-coherent algebras on $X$ and affine
	schemes over $X$.
	
	The cosimplicial diagram
	$\QCohAlg(U_\bullet)$ yields the cosimplicial diagram of affine schemes
	$\Aff/U_\bullet$. Since $\underline{\Spec}$ is compatible with base change,
	it preserves the descent data.
	
	Therefore the equivalence for quasi-coherent algebras induces an equivalence
	\[
	\Aff/X \cong \Tot\big(\Aff/U_\bullet\big).
	\]
\end{proof}

\begin{remark}
	Descent for affine schemes (and more generally affine morphisms) with
	respect to fpqc coverings has been prooven in the Stacks Project,
	see \cite[Tag 0244]{stacks-project}. The reduction to quasi-coherent
	algebras and the compatibility with relative $\Spec$ are standard and
	follow from fpqc descent for quasi-coherent sheaves,
	see \cite[Tag 023T]{stacks-project}.
\end{remark}


\begin{remark}
	This corollary is used repeatedly in quotient constructions: one first
	constructs the desired object locally on an fpqc cover, often after reducing
	to affine charts, and then glues the local affine pieces by fpqc descent.
	In particular, auxiliary affine morphisms arising in the proof of the
	Keel--Mori theorem may be constructed locally and descended globally.
\end{remark}